\chapter{Größen}

\startframedtext[width=\dimexpr\textwidth+\rightmarginwidth+\rightmargindistance\relax,frame=off,offset=none] % This spans textwidth + margin width

\startplacetable[location=here,title={Basisgrößen und Basiseinheiten in der Physik}]
\setupTABLE[each][align={middle,lohi}]
\setupTABLE[c][5][align={flushleft}]
\setupTABLE[row][1][background=color,backgroundcolor=gray,align={flushleft,lohi}]
\setupTABLE[row][2][background=color,backgroundcolor=gray]
\bTABLE[option=stretch,offset=1mm]
\bTR
	\bTD[nc=2] Basisgröße \eTD
	\bTD[nc=3] Basiseinheit \eTD
\eTR
\startCSV
Name, Symbol, Name, Symbol, Beschreibung
Zeit, $t$, Sekunde, $s$, "Die Sekunde, Synbol $s$, ist die die SI-Einheit der Zeit."
Länge, $l$, Meter, $m$, "Der Meter, Symbol $m$, ist die SI-Einheit der Länge."
Masse, $m$, Kilogramm, $kg$, "Das Kilogramm, Symbol $kg$, ist die SI-Einheit der Masse."
elektrische Stromstärke, $I$, Amper, $A$, "Das Ampere, Symbol $A$ , ist die SI-Einheit der elektrischen Stromstärke."
Temperatur, $T$, Kelvin, $K$, "Das Kelvin, Symbol $K$ , ist die SI-Einheit der thermodynamischen Temperatur."
Stoffmenge, $n$, Mol, $mol$, "Die Stoffmenge, Symbol $n$ , eines Systems ist ein Maß für die Anzahl spezifizierter elementarer Einzelteile."
Lichtstärke, $I_v$, Candela, $cd$, "Das Candela, Symbol $cd$, ist die SI-Einheit der Lichtstärke in einer vorgegebenen Richtung."
\stopCSV

\eTABLE
\stopplacetable

\stopframedtext % }}}

